
\documentclass{article}
\usepackage{colortbl}
\usepackage{makecell}
\usepackage{multirow}
\usepackage{supertabular}

\begin{document}

\newcounter{utterance}

\centering \large Interaction Transcript for game `dond', experiment `coop', episode 4 with gemma{-}4{-}26b{-}a4b{-}it.
\vspace{24pt}

{ \footnotesize  \setcounter{utterance}{1}
\setlength{\tabcolsep}{0pt}
\begin{supertabular}{c@{$\;$}|p{.15\linewidth}@{}p{.15\linewidth}p{.15\linewidth}p{.15\linewidth}p{.15\linewidth}p{.15\linewidth}}
    \# & $\;$A & \multicolumn{4}{c}{Game Master} & $\;\:$B\\
    \hline

    \theutterance \stepcounter{utterance}  
    & & \multicolumn{4}{p{0.6\linewidth}}{
        \cellcolor[rgb]{0.9,0.9,0.9}{
            \makecell[{{p{\linewidth}}}]{
                \texttt{\tiny{[P1$\langle$GM]}}
                \texttt{أنت تلعب لعبة تفاوض تعاونية يتعين عليك فيها الاتفاق على كيفية تقسيم مجموعة من العناصر بينك وبين لاعب آخر.} \\
\\ 
\texttt{القواعد:} \\
\texttt{(a) يتم إعطاؤك أنت واللاعب الآخر مجموعة من العناصر. كما يُعطى كل منكما دالة قيمة سرية، تمثل مقدار تقديرك لكل نوع من الأشياء.} \\
\texttt{(b) تتبادل الرسائل مع اللاعب الآخر للاتفاق على من يحصل على أي عناصر. يمكن لكل منكما إرسال حد أقصى قدره 5 رسالة، أو الإنهاء مبكرًا عبر تقديم اقتراح سري في أي وقت.} \\
\texttt{(c) يُطلب من كل منكما تقديم اقتراح سري يحدد العناصر التي تريدها بصيغة داخل أقواس مربعة كما يلي: "{[}اقتراح: <رقم> <اسم الكائن>, <رقم> <اسم الكائن>, <...>{]}"} \\
\texttt{(d) إذا كانت اقتراحاتكما متكاملة، أي توجد عناصر كافية لتلبية كلا الاقتراحين، فسيُمنح كل لاعب نتيجة بناءً على مجموع القيم للعناصر التي حصل عليها. وإلا، فسيحصل كلاكما على صفر نقاط.} \\
\texttt{(e) هدفك هو تعظيم مجموع نتيجتك ونتيجة اللاعب الآخر.} \\
\\ 
\texttt{لنبدأ.} \\
\\ 
\texttt{مجموعة العناصر المتاحة هي:} \\
\\ 
\texttt{1 كشتبان, 1 صابون, 4 طوابع} \\
\\ 
\texttt{دالة القيمة السرية الخاصة بك هي:} \\
\\ 
\texttt{كشتبان: 3, صابون: 7, طابع: 0} \\
\\ 
\texttt{مهم: رسائلك، ما لم تكن الاقتراح السري، تُنقل مباشرة إلى اللاعب الآخر، لذا لا تُضمّن أي رد على القواعد أو نص يعلن عن رسالتك. لتقديم اقتراح سري، استخدم الصيغة المشار إليها. لا تستخدم الأقواس المربعة عند التواصل مع اللاعب الآخر وإلا فسيُفسَّر ذلك على أنه اقتراحك السري.} \\
            }
        }
    }
    & \\ \\

    \theutterance \stepcounter{utterance}  
    & \multicolumn{4}{p{0.6\linewidth}}{
        \cellcolor[rgb]{0.9,0.9,0.9}{
            \makecell[{{p{\linewidth}}}]{
                \texttt{\tiny{[P1$\rangle$GM]}}
                \texttt{مرحباً. لدي مجموعة من العناصر وأرغب في الوصول إلى اتفاق يرضينا معاً. بما أن هدفي هو تعظيم مجموع نقاطنا، أقترح أن نوزع العناصر بناءً على ما يحتاجه كل منا.} \\
\\ 
\texttt{هل يمكنك إخباري ما هي العناصر التي تهمك أكثر؟ أو إذا كنت تفضل، يمكننا محاولة تقسيمها مباشرة.} \\
            }
        }
    }
    & & \\ \\

    \theutterance \stepcounter{utterance}  
    & & & \multicolumn{2}{p{0.3\linewidth}}{
        \cellcolor[rgb]{0.9,0.9,0.9}{
            \makecell[{{p{\linewidth}}}]{
                \texttt{\tiny{[GM$|$GM]}}
                \texttt{continue} \\
            }
        }
    }
    & & \\ \\

    \theutterance \stepcounter{utterance}  
    & & & \multicolumn{2}{p{0.3\linewidth}}{
        \cellcolor[rgb]{0.9,0.9,0.9}{
            \makecell[{{p{\linewidth}}}]{
                \texttt{\tiny{[GM$|$GM]}}
                \texttt{مرحباً. لدي مجموعة من العناصر وأرغب في الوصول إلى اتفاق يرضينا معاً. بما أن هدفي هو تعظيم مجموع نقاطنا، أقترح أن نوزع العناصر بناءً على ما يحتاجه كل منا.} \\
\\ 
\texttt{هل يمكنك إخباري ما هي العناصر التي تهمك أكثر؟ أو إذا كنت تفضل، يمكننا محاولة تقسيمها مباشرة.} \\
            }
        }
    }
    & & \\ \\

    \theutterance \stepcounter{utterance}  
    & & \multicolumn{4}{p{0.6\linewidth}}{
        \cellcolor[rgb]{0.9,0.9,0.9}{
            \makecell[{{p{\linewidth}}}]{
                \texttt{\tiny{[GM$\rangle$P2]}}
                \texttt{أنت تلعب لعبة تفاوض تعاونية يتعين عليك فيها الاتفاق على كيفية تقسيم مجموعة من العناصر بينك وبين لاعب آخر.} \\
\\ 
\texttt{القواعد:} \\
\texttt{(a) يتم إعطاؤك أنت واللاعب الآخر مجموعة من العناصر. كما يُعطى كل منكما دالة قيمة سرية، تمثل مقدار تقديرك لكل نوع من الأشياء.} \\
\texttt{(b) تتبادل الرسائل مع اللاعب الآخر للاتفاق على من يحصل على أي عناصر. يمكن لكل منكما إرسال حد أقصى قدره 5 رسالة، أو الإنهاء مبكرًا عبر تقديم اقتراح سري في أي وقت.} \\
\texttt{(c) يُطلب من كل منكما تقديم اقتراح سري يحدد العناصر التي تريدها بصيغة داخل أقواس مربعة كما يلي: "{[}اقتراح: <رقم> <اسم الكائن>, <رقم> <اسم الكائن>, <...>{]}"} \\
\texttt{(d) إذا كانت اقتراحاتكما متكاملة، أي توجد عناصر كافية لتلبية كلا الاقتراحين، فسيُمنح كل لاعب نتيجة بناءً على مجموع القيم للعناصر التي حصل عليها. وإلا، فسيحصل كلاكما على صفر نقاط.} \\
\texttt{(e) هدفك هو تعظيم مجموع نتيجتك ونتيجة اللاعب الآخر.} \\
\\ 
\texttt{لنبدأ.} \\
\\ 
\texttt{مجموعة العناصر المتاحة هي:} \\
\\ 
\texttt{1 كشتبان, 1 صابون, 4 طوابع} \\
\\ 
\texttt{دالة القيمة السرية الخاصة بك هي:} \\
\\ 
\texttt{كشتبان: 2, صابون: 0, طابع: 2} \\
\\ 
\texttt{مهم: رسائلك، ما لم تكن الاقتراح السري، تُنقل مباشرة إلى اللاعب الآخر، لذا لا تُضمّن أي رد على القواعد أو نص يعلن عن رسالتك. لتقديم اقتراح سري، استخدم الصيغة المشار إليها. لا تستخدم الأقواس المربعة عند التواصل مع اللاعب الآخر وإلا فسيُفسَّر ذلك على أنه اقتراحك السري.} \\
\\ 
\texttt{مرحباً. لدي مجموعة من العناصر وأرغب في الوصول إلى اتفاق يرضينا معاً. بما أن هدفي هو تعظيم مجموع نقاطنا، أقترح أن نوزع العناصر بناءً على ما يحتاجه كل منا.} \\
\\ 
\texttt{هل يمكنك إخباري ما هي العناصر التي تهمك أكثر؟ أو إذا كنت تفضل، يمكننا محاولة تقسيمها مباشرة.} \\
            }
        }
    }
    & \\ \\

    \theutterance \stepcounter{utterance}  
    & & & \multicolumn{4}{p{0.6\linewidth}}{
        \cellcolor[rgb]{0.9,0.9,0.9}{
            \makecell[{{p{\linewidth}}}]{
                \texttt{\tiny{[GM$\langle$P2]}}
                \texttt{أهلاً بك. أنا أوافقك الرأي تماماً، فهدفنا المشترك هو تعظيم المنفعة الإجمالية للطرفين.} \\
\\ 
\texttt{بما أن العناصر محدودة (كشتبان واحد، صابون واحد، و4 طوابع)، فإن المفتاح هو تجنب التكرار في الطلبات لضمان عدم حصولنا على صفر.} \\
\\ 
\texttt{أنا أهتم بشكل أساسي بالكشتبان والطوابع، بينما لا أرى قيمة في الصابون. هل يمكنك إخباري ما هي العناصر التي تمنحك أعلى قيمة؟ هذا سيساعدنا على تقسيمها بدقة دون تداخل.} \\
            }
        }
    }
     \\ \\

    \theutterance \stepcounter{utterance}  
    & & & \multicolumn{2}{p{0.3\linewidth}}{
        \cellcolor[rgb]{0.9,0.9,0.9}{
            \makecell[{{p{\linewidth}}}]{
                \texttt{\tiny{[GM$|$GM]}}
                \texttt{continue} \\
            }
        }
    }
    & & \\ \\

    \theutterance \stepcounter{utterance}  
    & & & \multicolumn{2}{p{0.3\linewidth}}{
        \cellcolor[rgb]{0.9,0.9,0.9}{
            \makecell[{{p{\linewidth}}}]{
                \texttt{\tiny{[GM$|$GM]}}
                \texttt{أهلاً بك. أنا أوافقك الرأي تماماً، فهدفنا المشترك هو تعظيم المنفعة الإجمالية للطرفين.} \\
\\ 
\texttt{بما أن العناصر محدودة (كشتبان واحد، صابون واحد، و4 طوابع)، فإن المفتاح هو تجنب التكرار في الطلبات لضمان عدم حصولنا على صفر.} \\
\\ 
\texttt{أنا أهتم بشكل أساسي بالكشتبان والطوابع، بينما لا أرى قيمة في الصابون. هل يمكنك إخباري ما هي العناصر التي تمنحك أعلى قيمة؟ هذا سيساعدنا على تقسيمها بدقة دون تداخل.} \\
            }
        }
    }
    & & \\ \\

    \theutterance \stepcounter{utterance}  
    & & \multicolumn{4}{p{0.6\linewidth}}{
        \cellcolor[rgb]{0.9,0.9,0.9}{
            \makecell[{{p{\linewidth}}}]{
                \texttt{\tiny{[P1$\langle$GM]}}
                \texttt{أهلاً بك. أنا أوافقك الرأي تماماً، فهدفنا المشترك هو تعظيم المنفعة الإجمالية للطرفين.} \\
\\ 
\texttt{بما أن العناصر محدودة (كشتبان واحد، صابون واحد، و4 طوابع)، فإن المفتاح هو تجنب التكرار في الطلبات لضمان عدم حصولنا على صفر.} \\
\\ 
\texttt{أنا أهتم بشكل أساسي بالكشتبان والطوابع، بينما لا أرى قيمة في الصابون. هل يمكنك إخباري ما هي العناصر التي تمنحك أعلى قيمة؟ هذا سيساعدنا على تقسيمها بدقة دون تداخل.} \\
            }
        }
    }
    & \\ \\

    \theutterance \stepcounter{utterance}  
    & \multicolumn{4}{p{0.6\linewidth}}{
        \cellcolor[rgb]{0.9,0.9,0.9}{
            \makecell[{{p{\linewidth}}}]{
                \texttt{\tiny{[P1$\rangle$GM]}}
                \texttt{شكراً لصراحتك، هذا يسهل الأمر كثيراً.} \\
\\ 
\texttt{بما أنك مهتم بالكشتبان والطوابع، وأنا أرى قيمة عالية جداً في الصابون (7 نقاط) بينما لا أهتم بالطوابع إطلاقاً (0 نقطة)، أعتقد أن الحل الأمثل هو أن آخذ الصابون لنفسي، وأن تذهب الكشتبان والطوابع الأربعة إليك بالكامل.} \\
\\ 
\texttt{بهذه الطريقة، نضمن عدم التداخل ونحقق أقصى استفادة ممكنة من العناصر المتاحة. ما رأيك في هذا التقسيم؟} \\
            }
        }
    }
    & & \\ \\

    \theutterance \stepcounter{utterance}  
    & & & \multicolumn{2}{p{0.3\linewidth}}{
        \cellcolor[rgb]{0.9,0.9,0.9}{
            \makecell[{{p{\linewidth}}}]{
                \texttt{\tiny{[GM$|$GM]}}
                \texttt{continue} \\
            }
        }
    }
    & & \\ \\

    \theutterance \stepcounter{utterance}  
    & & & \multicolumn{2}{p{0.3\linewidth}}{
        \cellcolor[rgb]{0.9,0.9,0.9}{
            \makecell[{{p{\linewidth}}}]{
                \texttt{\tiny{[GM$|$GM]}}
                \texttt{شكراً لصراحتك، هذا يسهل الأمر كثيراً.} \\
\\ 
\texttt{بما أنك مهتم بالكشتبان والطوابع، وأنا أرى قيمة عالية جداً في الصابون (7 نقاط) بينما لا أهتم بالطوابع إطلاقاً (0 نقطة)، أعتقد أن الحل الأمثل هو أن آخذ الصابون لنفسي، وأن تذهب الكشتبان والطوابع الأربعة إليك بالكامل.} \\
\\ 
\texttt{بهذه الطريقة، نضمن عدم التداخل ونحقق أقصى استفادة ممكنة من العناصر المتاحة. ما رأيك في هذا التقسيم؟} \\
            }
        }
    }
    & & \\ \\

    \theutterance \stepcounter{utterance}  
    & & \multicolumn{4}{p{0.6\linewidth}}{
        \cellcolor[rgb]{0.9,0.9,0.9}{
            \makecell[{{p{\linewidth}}}]{
                \texttt{\tiny{[GM$\rangle$P2]}}
                \texttt{شكراً لصراحتك، هذا يسهل الأمر كثيراً.} \\
\\ 
\texttt{بما أنك مهتم بالكشتبان والطوابع، وأنا أرى قيمة عالية جداً في الصابون (7 نقاط) بينما لا أهتم بالطوابع إطلاقاً (0 نقطة)، أعتقد أن الحل الأمثل هو أن آخذ الصابون لنفسي، وأن تذهب الكشتبان والطوابع الأربعة إليك بالكامل.} \\
\\ 
\texttt{بهذه الطريقة، نضمن عدم التداخل ونحقق أقصى استفادة ممكنة من العناصر المتاحة. ما رأيك في هذا التقسيم؟} \\
            }
        }
    }
    & \\ \\

    \theutterance \stepcounter{utterance}  
    & & & \multicolumn{4}{p{0.6\linewidth}}{
        \cellcolor[rgb]{0.9,0.9,0.9}{
            \makecell[{{p{\linewidth}}}]{
                \texttt{\tiny{[GM$\langle$P2]}}
                \texttt{هذا اقتراح ممتاز ومنطقي جداً! بما أنك تمنح الصابون قيمة عالية جداً (7 نقاط) وأنا لا أراه ذا قيمة (0 نقطة)، بينما أنا أهتم بالطوابع وأنت لا تهتم بها، فإن هذا التقسيم يحقق أقصى مجموع ممكن للنقاط لكلينا دون أي تداخل.} \\
\\ 
\texttt{بناءً على ذلك، سأقدم اقتراحي السري الآن لضمان تنفيذ هذا الاتفاق.} \\
            }
        }
    }
     \\ \\

    \theutterance \stepcounter{utterance}  
    & & & \multicolumn{2}{p{0.3\linewidth}}{
        \cellcolor[rgb]{0.9,0.9,0.9}{
            \makecell[{{p{\linewidth}}}]{
                \texttt{\tiny{[GM$|$GM]}}
                \texttt{continue} \\
            }
        }
    }
    & & \\ \\

    \theutterance \stepcounter{utterance}  
    & & & \multicolumn{2}{p{0.3\linewidth}}{
        \cellcolor[rgb]{0.9,0.9,0.9}{
            \makecell[{{p{\linewidth}}}]{
                \texttt{\tiny{[GM$|$GM]}}
                \texttt{هذا اقتراح ممتاز ومنطقي جداً! بما أنك تمنح الصابون قيمة عالية جداً (7 نقاط) وأنا لا أراه ذا قيمة (0 نقطة)، بينما أنا أهتم بالطوابع وأنت لا تهتم بها، فإن هذا التقسيم يحقق أقصى مجموع ممكن للنقاط لكلينا دون أي تداخل.} \\
\\ 
\texttt{بناءً على ذلك، سأقدم اقتراحي السري الآن لضمان تنفيذ هذا الاتفاق.} \\
            }
        }
    }
    & & \\ \\

    \theutterance \stepcounter{utterance}  
    & & \multicolumn{4}{p{0.6\linewidth}}{
        \cellcolor[rgb]{0.9,0.9,0.9}{
            \makecell[{{p{\linewidth}}}]{
                \texttt{\tiny{[P1$\langle$GM]}}
                \texttt{هذا اقتراح ممتاز ومنطقي جداً! بما أنك تمنح الصابون قيمة عالية جداً (7 نقاط) وأنا لا أراه ذا قيمة (0 نقطة)، بينما أنا أهتم بالطوابع وأنت لا تهتم بها، فإن هذا التقسيم يحقق أقصى مجموع ممكن للنقاط لكلينا دون أي تداخل.} \\
\\ 
\texttt{بناءً على ذلك، سأقدم اقتراحي السري الآن لضمان تنفيذ هذا الاتفاق.} \\
            }
        }
    }
    & \\ \\

    \theutterance \stepcounter{utterance}  
    & \multicolumn{4}{p{0.6\linewidth}}{
        \cellcolor[rgb]{0.9,0.9,0.9}{
            \makecell[{{p{\linewidth}}}]{
                \texttt{\tiny{[P1$\rangle$GM]}}
                \texttt{{[}اقتراح: 1 صابون{]}} \\
            }
        }
    }
    & & \\ \\

    \theutterance \stepcounter{utterance}  
    & & & \multicolumn{2}{p{0.3\linewidth}}{
        \cellcolor[rgb]{0.9,0.9,0.9}{
            \makecell[{{p{\linewidth}}}]{
                \texttt{\tiny{[GM$|$GM]}}
                \texttt{proposal} \\
            }
        }
    }
    & & \\ \\

    \theutterance \stepcounter{utterance}  
    & & & \multicolumn{2}{p{0.3\linewidth}}{
        \cellcolor[rgb]{0.9,0.9,0.9}{
            \makecell[{{p{\linewidth}}}]{
                \texttt{\tiny{[GM$|$GM]}}
                [0, 1, 0]
            }
        }
    }
    & & \\ \\

    \theutterance \stepcounter{utterance}  
    & & \multicolumn{4}{p{0.6\linewidth}}{
        \cellcolor[rgb]{0.9,0.9,0.9}{
            \makecell[{{p{\linewidth}}}]{
                \texttt{\tiny{[GM$\rangle$P2]}}
                \texttt{قدّم اللاعب الآخر اقتراحًا سريًا. حان الآن دورك لتقديم اقتراحك الخاص.} \\
            }
        }
    }
    & \\ \\

    \theutterance \stepcounter{utterance}  
    & & & \multicolumn{4}{p{0.6\linewidth}}{
        \cellcolor[rgb]{0.9,0.9,0.9}{
            \makecell[{{p{\linewidth}}}]{
                \texttt{\tiny{[GM$\langle$P2]}}
                \texttt{{[}اقتراح: 1 كشتبان, 4 طابع{]}} \\
            }
        }
    }
     \\ \\

    \theutterance \stepcounter{utterance}  
    & & & \multicolumn{2}{p{0.3\linewidth}}{
        \cellcolor[rgb]{0.9,0.9,0.9}{
            \makecell[{{p{\linewidth}}}]{
                \texttt{\tiny{[GM$|$GM]}}
                \texttt{proposal} \\
            }
        }
    }
    & & \\ \\

    \theutterance \stepcounter{utterance}  
    & & & \multicolumn{2}{p{0.3\linewidth}}{
        \cellcolor[rgb]{0.9,0.9,0.9}{
            \makecell[{{p{\linewidth}}}]{
                \texttt{\tiny{[GM$|$GM]}}
                [1, 0, 4]
            }
        }
    }
    & & \\ \\

    \theutterance \stepcounter{utterance}  
    & & & \multicolumn{2}{p{0.3\linewidth}}{
        \cellcolor[rgb]{0.9,0.9,0.9}{
            \makecell[{{p{\linewidth}}}]{
                \texttt{\tiny{[GM$|$GM]}}
                [[0, 1, 0], [1, 0, 4]]
            }
        }
    }
    & & \\ \\

\end{supertabular}
}

\end{document}
